\documentclass{article}
\usepackage{amsmath,amssymb,amsthm}
\usepackage{algorithm}
\usepackage{algpseudocode}
\usepackage{fullpage}
\newtheorem{theorem}{Theorem}
\newtheorem{lemma}{Lemma}
\newtheorem{assumption}{Assumption}
\newtheorem{corollary}{Corollary}
\begin{document}

\section*{Random Block Coordinate Gradient Descent (RBCGD)}

\section*{Mathematical Formulation}
Let $f:\mathbb{R}^{d}\rightarrow\mathbb{R}$ be differentiable and possibly non‑convex.  
Fix a partition of the coordinate set $\{1,\dots ,d\}$ into $B$ disjoint blocks
\[
\mathcal{B}_{1},\dots ,\mathcal{B}_{B},\qquad 
\bigcup_{i=1}^{B}\mathcal{B}_{i}=\{1,\dots ,d\},\; \mathcal{B}_{i}\cap\mathcal{B}_{j}=\emptyset\;(i\neq j).
\]
Denote by $w_{\mathcal{B}_{i}}$ the sub‑vector of $w$ restricted to block $\mathcal{B}_{i}$ and by
$\nabla_{\mathcal{B}_{i}}f(w)$ the corresponding partial gradient.

The algorithm draws at each iteration $t\ge 0$ a block index $i_{t}\in\{1,\dots ,B\}$ uniformly at random,
computes the partial gradient $\nabla_{\mathcal{B}_{i_{t}}}f(w_{t})$ and updates only that block:
\[
\boxed{\,w_{t+1}= w_{t} - \eta_{t}\,U_{i_{t}}\nabla_{\mathcal{B}_{i_{t}}}f(w_{t})\,},
\tag{1}
\]
where $\eta_{t}>0$ is a stepsize and $U_{i}$ is the $d\times |\mathcal{B}_{i}|$ zero‑padding operator
\[
U_{i}v=
\begin{bmatrix}
0\\ \vdots \\ v \\ \vdots \\ 0
\end{bmatrix},
\qquad
\text{the vector $v$ is placed on the coordinates of block }\mathcal{B}_{i}.
\]

\section*{Pseudocode}
\begin{algorithm}[H]
\caption{Random Block Coordinate Gradient Descent}
\begin{algorithmic}[1]
\State \textbf{Input:} initial point $w_{0}\in\mathbb{R}^{d}$, stepsize schedule $\{\eta_{t}\}_{t\ge0}$, number of iterations $T$.
\For{$t=0$ \textbf{to} $T-1$}
    \State Sample $i_{t}\sim\text{Uniform}\{1,\dots ,B\}$.
    \State Compute partial gradient $g_{t}= \nabla_{\mathcal{B}_{i_{t}}}f(w_{t})$.
    \State $w_{t+1}= w_{t} - \eta_{t}\,U_{i_{t}} g_{t}$.
\EndFor
\State \textbf{Output:} $w_{T}$ (or the average $\bar w_{T}= \frac1T\sum_{t=0}^{T-1} w_{t}$).
\end{algorithmic}
\end{algorithm}

\section*{Dry Run Example}
Consider the scalar function $f(w)=(w-3)^{2}$, i.e. $d=1$, a single block $B=1$, and stepsize $\eta=0.1$.
\[
\nabla f(w)=2(w-3).
\]

\begin{center}
\begin{tabular}{c|c|c|c}
$t$ & $w_{t}$ & $\nabla f(w_{t})$ & $f(w_{t})$ \\ \hline
0 & $0$ & $-6$ & $9$ \\
1 & $w_{1}=0-0.1(-6)=0.6$ & $2(0.6-3)=-4.8$ & $(0.6-3)^{2}=5.76$ \\
2 & $w_{2}=0.6-0.1(-4.8)=1.08$ & $2(1.08-3)=-3.84$ & $(1.08-3)^{2}=3.6864$ \\
3 & $w_{3}=1.08-0.1(-3.84)=1.464$ & $2(1.464-3)=-3.072$ & $(1.464-3)^{2}=2.3630$
\end{tabular}
\end{center}
The objective value decreases at each iteration, illustrating the basic behaviour of (1).

\newpage
\section*{Assumptions}
\begin{assumption}[Blockwise $L$‑smoothness]\label{ass:Lsmooth}
For each block $i\in\{1,\dots ,B\}$ there exists $L_{i}>0$ such that for all $w,\tilde w\in\mathbb{R}^{d}$
\[
f\bigl(w+U_{i}( \tilde w_{\mathcal{B}_{i}}-w_{\mathcal{B}_{i}})\bigr)
\le f(w) + \langle \nabla_{\mathcal{B}_{i}}f(w),\tilde w_{\mathcal{B}_{i}}-w_{\mathcal{B}_{i}}\rangle
   + \frac{L_{i}}{2}\,\bigl\|\tilde w_{\mathcal{B}_{i}}-w_{\mathcal{B}_{i}}\bigr\|^{2}.
\]
\end{assumption}
Define $L_{\max}:=\max_{i}L_{i}$.

\begin{assumption}[Lower boundedness]\label{ass:lower}
$f$ is bounded from below: $\exists f^{\inf}\in\mathbb{R}$ such that $f(w)\ge f^{\inf}$ for all $w$.
\end{assumption}

\begin{assumption}[Stepsize]\label{ass:stepsize}
The stepsizes satisfy $0<\eta_{t}\le \frac{1}{L_{\max}}$ for all $t$.
\end{assumption}

\section*{Main Convergence Theorem}
\begin{theorem}\label{thm:main}
Assume \ref{ass:Lsmooth}--\ref{ass:stepsize}. Let $\{w_{t}\}$ be generated by RBCGD (1) and define the random block gradient norm
\[
G_{t}:=\bigl\|\nabla_{\mathcal{B}_{i_{t}}}f(w_{t})\bigr\|^{2}.
\]
Then for any $T\ge1$,
\[
\frac{1}{T}\sum_{t=0}^{T-1}\mathbb{E}\!\left[G_{t}\right]
\;\le\;
\frac{2L_{\max}\bigl(f(w_{0})-f^{\inf}\bigr)}{T}. \tag{2}
\]
Consequently,
\[
\min_{0\le t<T}\mathbb{E}\!\left[\bigl\|\nabla f(w_{t})\bigr\|^{2}\right]
\;\le\;
\frac{2BL_{\max}\bigl(f(w_{0})-f^{\inf}\bigr)}{T}. \tag{3}
\]
Thus the method attains the standard $O(1/T)$ rate for the expected squared gradient norm in the non‑convex setting.
\end{theorem}

\section*{Theoretical Properties}
\begin{itemize}
\item \textbf{Stability.} Because $\eta_{t}\le 1/L_{\max}$, (1) is a non‑expansive map in expectation; the iterates remain in a sublevel set of $f$ and are therefore bounded.
\item \textbf{Hyperparameter sensitivity.} The only required tuning is an upper bound on $L_{\max}$. Choosing $\eta$ smaller than $1/L_{\max}$ only slows convergence (the bound (2) scales linearly with $\eta$), while a larger $\eta$ may violate the descent Lemma 1 and break the guarantee.
\item \textbf{Comparison to full‑gradient GD.} For the same stepsize $\eta\le 1/L_{\max}$ full gradient descent yields
$\frac{1}{T}\sum_{t}\mathbb{E}\|\nabla f(w_{t})\|^{2}\le \frac{2L_{\max}(f(w_{0})-f^{\inf})}{T}$.
RBCGD incurs an extra factor $B$ in (3) because only one block is updated per iteration; however each iteration is $1/B$ of the computational cost of a full gradient step.
\end{itemize}

\section*{Discussion}
The theorem shows that even without convexity the random block scheme drives the expected squared norm of the gradient to zero at the optimal $O(1/T)$ rate. The proof hinges on the blockwise smoothness (Assumption~\ref{ass:Lsmooth}) and on the unbiasedness of the random block selection. Extensions to non‑uniform sampling, variance‑reduced estimators, or adaptive stepsizes follow by minor modifications of Lemma~1 and the telescoping sum used in the proof.

\newpage
\section*{Preliminaries (Definitions and Lemmas)}
\begin{definition}[Full gradient norm in terms of blocks]
For any $w$, because the blocks form a partition,
\[
\|\nabla f(w)\|^{2}= \sum_{i=1}^{B}\bigl\|\nabla_{\mathcal{B}_{i}}f(w)\bigr\|^{2}. \tag{4}
\]
\end{definition}

\begin{lemma}[Block Descent Lemma]\label{lem:descent}
Let $i\in\{1,\dots ,B\}$ and $w'\!=\!w-U_{i}\eta\nabla_{\mathcal{B}_{i}}f(w)$ with $0<\eta\le 1/L_{i}$. Then
\[
f(w')\le f(w)-\frac{\eta}{2}\,\bigl\|\nabla_{\mathcal{B}_{i}}f(w)\bigr\|^{2}. \tag{5}
\]
\end{lemma}
\begin{proof}
By Assumption~\ref{ass:Lsmooth} with $\tilde w_{\mathcal{B}_{i}}=w_{\mathcal{B}_{i}}-\eta\nabla_{\mathcal{B}_{i}}f(w)$,
\[
\begin{aligned}
f(w') &\le f(w)-\eta\bigl\|\nabla_{\mathcal{B}_{i}}f(w)\bigr\|^{2}
        +\frac{L_{i}\eta^{2}}{2}\bigl\|\nabla_{\mathcal{B}_{i}}f(w)\bigr\|^{2}}\\
      &= f(w)-\eta\Bigl(1-\frac{L_{i}\eta}{2}\Bigr)\bigl\|\nabla_{\mathcal{B}_{i}}f(w)\bigr\|^{2}.
\end{aligned}
\]
Since $\eta\le 1/L_{i}$, the factor in parentheses is at least $1/2$, yielding (5). \hfill\qedsymbol
\end{proof}

\section*{Main Proof (Step‑by‑Step Derivation)}
We prove Theorem~\ref{thm:main}. Throughout the proof, expectations are taken over the randomness of the block index $i_{t}$ conditioned on the history up to time $t$.

\noindent\textbf{[Step 1]} Apply Lemma~\ref{lem:descent} conditionally on $i_{t}$:
\[
\mathbb{E}\!\bigl[f(w_{t+1})\mid w_{t}\bigr]
\le f(w_{t})-\frac{\eta_{t}}{2}\,
      \mathbb{E}\!\bigl[\bigl\|\nabla_{\mathcal{B}_{i_{t}}}f(w_{t})\bigr\|^{2}\mid w_{t}\bigr].
\tag{6}
\]

\noindent\textbf{[Step 2]} Because $i_{t}$ is drawn uniformly,
\[
\mathbb{E}\!\bigl[\bigl\|\nabla_{\mathcal{B}_{i_{t}}}f(w_{t})\bigr\|^{2}\mid w_{t}\bigr]
   =\frac{1}{B}\sum_{i=1}^{B}\bigl\|\nabla_{\mathcal{B}_{i}}f(w_{t})\bigr\|^{2}
   =\frac{1}{B}\,\bigl\|\nabla f(w_{t})\bigr\|^{2} \quad\text{by (4)}.
\tag{7}
\]

\noindent\textbf{[Step 3]} Substitute (7) into (6):
\[
\mathbb{E}\!\bigl[f(w_{t+1})\mid w_{t}\bigr]
\le f(w_{t})-\frac{\eta_{t}}{2B}\,\bigl\|\nabla f(w_{t})\bigr\|^{2}.
\tag{8}
\]

\noindent\textbf{[Step 4]} Take total expectation and use the law of total expectation:
\[
\mathbb{E}\bigl[f(w_{t+1})\bigr]
\le \mathbb{E}\bigl[f(w_{t})\bigr]-\frac{\eta_{t}}{2B}\,
      \mathbb{E}\bigl[\bigl\|\nabla f(w_{t})\bigr\|^{2}\bigr].
\tag{9}
\]

\noindent\textbf{[Step 5]} Rearrange (9) to isolate the gradient term:
\[
\frac{\eta_{t}}{2B}\,
\mathbb{E}\bigl[\bigl\|\nabla f(w_{t})\bigr\|^{2}\bigr]
\le \mathbb{E}\bigl[f(w_{t})\bigr]-\mathbb{E}\bigl[f(w_{t+1})\bigr].
\tag{10}
\]

\noindent\textbf{[Step 6]} Sum (10) over $t=0,\dots ,T-1$ and telescope:
\[
\sum_{t=0}^{T-1}\frac{\eta_{t}}{2B}\,
\mathbb{E}\bigl[\bigl\|\nabla f(w_{t})\bigr\|^{2}\bigr]
\le \mathbb{E}\bigl[f(w_{0})\bigr]-\mathbb{E}\bigl[f(w_{T})\bigr]
\le f(w_{0})-f^{\inf},
\tag{11}
\]
where the last inequality uses Assumption~\ref{ass:lower}.

\noindent\textbf{[Step 7]} Choose a *constant* stepsize $\eta_{t}\equiv\eta\le 1/L_{\max}$.
Then (11) becomes
\[
\frac{\eta}{2B}\sum_{t=0}^{T-1}
\mathbb{E}\bigl[\bigl\|\nabla f(w_{t})\bigr\|^{2}\bigr]
\le f(w_{0})-f^{\inf}.
\tag{12}
\]

\noindent\textbf{[Step 8]} Divide both sides of (12) by $T$ and multiply by $2B/\eta$:
\[
\frac{1}{T}\sum_{t=0}^{T-1}
\mathbb{E}\bigl[\bigl\|\nabla f(w_{t})\bigr\|^{2}\bigr]
\le \frac{2B\bigl(f(w_{0})-f^{\inf}\bigr)}{\eta T}.
\tag{13}
\]

\noindent\textbf{[Step 9]} Since $\eta\le 1/L_{\max}$, we have $1/\eta\le L_{\max}$. Plugging this bound into (13) yields
\[
\frac{1}{T}\sum_{t=0}^{T-1}
\mathbb{E}\bigl[\bigl\|\nabla f(w_{t})\bigr\|^{2}\bigr]
\le \frac{2BL_{\max}\bigl(f(w_{0})-f^{\inf}\bigr)}{T}.
\tag{14}
\]

\noindent\textbf{[Step 10]} Relation (14) immediately implies the existence of an iterate with small expected gradient norm:
\[
\min_{0\le t<T}
\mathbb{E}\bigl[\bigl\|\nabla f(w_{t})\bigr\|^{2}\bigr]
\le
\frac{2BL_{\max}\bigl(f(w_{0})-f^{\inf}\bigr)}{T},
\tag{15}
\]
which is exactly the bound (3) of Theorem~\ref{thm:main}.

\noindent\textbf{[Step 11]} To obtain the bound (2) on the *block* gradient norm $G_{t}$, note from (7) that
\[
\mathbb{E}[G_{t}]
= \frac{1}{B}\,\mathbb{E}\bigl[\|\nabla f(w_{t})\|^{2}\bigr].
\]
Insert this relation into (14) and simplify:
\[
\frac{1}{T}\sum_{t=0}^{T-1}\mathbb{E}[G_{t}]
\le
\frac{2L_{\max}\bigl(f(w_{0})-f^{\inf}\bigr)}{T},
\]
which is precisely (2). ∎

\section*{Rate Derivation (With Constants)}
The constants appearing in the final bound are:
\[
C_{1}=2L_{\max}\bigl(f(w_{0})-f^{\inf}\bigr),\qquad
C_{2}=2BL_{\max}\bigl(f(w_{0})-f^{\inf}\bigr).
\]
Thus the expected block gradient norm decays as $C_{1}/T$ and the full gradient norm as $C_{2}/T$.

\section*{Special Cases}
\begin{itemize}
\item \textbf{Uniform Lipschitz constants.} If $L_{i}=L$ for all $i$, then $L_{\max}=L$ and the bound simplifies to $2L(f(w_{0})-f^{\inf})/T$ for the block gradient.
\item \textbf{Non‑uniform sampling.} Let block $i$ be selected with probability $p_{i}>0$, $\sum_{i}p_{i}=1$, and define $L_{\max}^{p}=\max_{i}L_{i}/p_{i}$. Repeating the proof with $p_{i}$ in place of $1/B$ yields the same $O(1/T)$ rate with $L_{\max}^{p}$ in place of $L_{\max}$.
\item \textbf{Decreasing stepsize.} If $\eta_{t}= \frac{c}{t+1}$ with $c\le 1/L_{\max}$, the telescoping argument gives
\[
\frac{1}{\sum_{t=0}^{T-1}\eta_{t}}\sum_{t=0}^{T-1}\eta_{t}\,
\mathbb{E}\bigl[\|\nabla f(w_{t})\|^{2}\bigr]
\le \frac{2B\bigl(f(w_{0})-f^{\inf}\bigr)}{\sum_{t=0}^{T-1}\eta_{t}},
\]
which is also $O(1/\log T)$; the constant stepsize choice is therefore optimal for the $O(1/T)$ rate.
\end{itemize}

\end{document}